\scanpage{596}

\mspage{5}

\deleted{Soit $\mathrm{Max}(k_E)$ l'ens \deleted{des} le spectre maximal de $k_E$,}

Pour $\unsure{\forall|\mathrm{tt}}$ anneau $A$, soit $\mathrm{Max}(A)$ \unsure{son|le} spectre maximal, i.e. l'ens des idéaux \unsure{maxima|maximaux} de $A$. On \unsure{trouve} donc une application canonique

\[
\left|\begin{aligned}
\unsure{\mathscr F|\mathcal F}_0&\longrightarrow\mathrm{Max}(k_E)\\
\Phi&\longmapsto\mathrm{Ker}(k_E\longrightarrow k_\Phi)\overset{\mathrm{\unsure{déf|def}}}{=}\unsure{\mathfrak m|m}_\Phi
\end{aligned}\right.\tag*{\makebox[0pt][r]{\makebox[\displaywidth][l]{\unsure{(a)|(A)}}}}
\]

\uline{Proposition 2}\qquad L'application précédente est \uline{bijective}.

\unsure{Prouvons} d'abord qu'elle est injective, cela \deleted{provient de} résulte aussitôt du

% slanted
\margin{Cf \unsure{ci-après} pour \unsure{le} cor. 3 c), pour $\Phi$ filtre quelconque}

\uline{Corollaire 1}\qquad Soit, pour $\Phi\in\mathscr F_0$,

\[
\unsure{\chi|X}_\Phi:k_E\mathrel{\unsure{-|\longrightarrow}}k_\Phi
\]

l'application canonique. Alors pour $A\in\unsure{\mathcal P|\mathfrak P}(E)$, on a

% slanted
\margin{\uline{NB} Pour $\Phi\in\mathscr F$ quelconque, $\mathfrak m_\Phi\overset{\mathrm{def}}{=}\mathrm{Ker}(\chi_\Phi:k_E\mathrel{\unsure{-|\to}} k_\Phi)$ détermine $\Phi$, puisque $A\in\Phi\Longleftrightarrow\chi_\Phi(\unsure{f_A|\varphi_A})=1$ (mais il faut \unsure{supposer} les $k_{\unsure{\xi|\zeta}}\ne0$, i.e. $0_\xi\ne1_\xi$ \unsure{\ldots})}

% slanted
\margin{\deleted{$k_\xi$} \unsure{remonter|remonté|anneaux} \uline{\unsure{avant|unit}} \unsure{corps}}

\[
\begin{aligned}
 A\notin\Phi&\Longleftrightarrow\chi_\Phi(f_A)\unsure{=|\ne}0\quad\text{i.e. } f_A\in\mathfrak m_\Phi\\
 \text{(i.e. }\ A\in\Phi&\Longleftrightarrow\chi_\Phi(f_A)\ne0\quad\text{i.e. } f_A\notin\mathfrak m_\Phi).
\end{aligned}
\]

% slanted
\margin{\uline{NB} \unsure{bien sûr} $(\chi_\Phi(f_A)\ne0)\Longleftrightarrow(\chi_\Phi(f_A)\unsure{=1|\ne1})$}

En effet, $\chi_\Phi(f_A)=0$ signifie qu'il existe $B\in\Phi$ tel que $f_A\mid B=0$, i.e. tel que $B\subset E\setminus A$; mais \deleted{\unsure{cela}} \deleted{$E\setminus A\in\Phi$, et} équivaut à $E\setminus A\in\Phi$, donc \unsure{à|:} $A\notin\Phi$ puisque $\Phi$ est un ultrafiltre, qed.

% archive sheet number: 595
